% Kratki čitateljski prijelazi za tiskano izdanje.
\ExplSyntaxOn
\NewDocumentCommand{\BookPartContains}{mmm}{\str_if_in:nnT{#1}{#2}{#3}}
\ExplSyntaxOff
\newcommand{\BookPartIntro}[1]{%
  \BookPartContains{#1}{DIO I:}{%
  Prije prvoga računa razjašnjavamo što broj predstavlja i koju tvrdnju
  može poduprijeti. Pratimo nastanak podataka, mjerenje i dizajn istraživanja,
  a zatim provjeravamo kako se točan broj može pretvoriti u pogrešan zaključak.
  }%
  \BookPartContains{#1}{DIO II:}{%
  Podatke sada sažimamo sredinama, mjerama raspršenosti i grafovima.
  Svaki prikaz nešto čuva, a nešto izostavlja. Učimo odabrati sažetak prema
  pitanju i opisati povezanost bez preuranjenoga uzročnog tumačenja.
  }%
  \BookPartContains{#1}{DIO III:}{%
  Opis uzorka nije još tvrdnja o populaciji. Kroz slučajnost i ponovljeno
  uzorkovanje pratimo kako se procjene mijenjaju. Iz te promjenjivosti
  gradimo intervale, uz granice koje određuju odabir i mjerenje.
  }%
  \BookPartContains{#1}{DIO IV:}{%
  Procjeni i intervalu pridružujemo usporedbu s nultim modelom.
  Proučavamo pogreške odluke, veličinu učinka i planiranje istraživanja,
  a zatim provjeravamo kako izbor analiza i objavljivanje oblikuju dokaze.
  }%
  \BookPartContains{#1}{DIO V:}{%
  Ista pitanja o procjeni, nesigurnosti i dosegu primjenjujemo na tablice
  kategorija, usporedbe skupina i regresiju. Modelom povezujemo više
  obilježja, a pri predviđanju provjeravamo kako model radi na novim podacima i
  kakve posljedice odluke imaju za ljude.
  }%
  \BookPartContains{#1}{ZAVRŠNICA}{%
  Povezujemo pitanje, podatke, model i izvještaj u jednu provjerljivu cjelinu.
  Pratimo jedan regresijski primjer i bilježimo odluke koje određuju njegov
  doseg. Asistent pomaže u radu, dok provjera i odgovornost ostaju naše.
  }%
}
